% Results
\section{Results}
\label{sec:results}

We report results for two existence hypotheses: user learning (reformulation slope) and AI responsiveness (diversity correlation). Both hypotheses are confirmed with statistical significance, though effect sizes are smaller than initially targeted. We also present unexpected findings regarding heterogeneity in learning patterns and temporal dynamics of responsiveness.

\subsection{Hypothesis 1: User Learning via Reformulation Slope}

\textbf{Research question}: Do users reformulate queries less frequently as conversations progress?

\textbf{Finding}: Yes. Mean reformulation slope is significantly negative, indicating users learn AI capabilities and reformulate less over turns.

\subsubsection{Aggregate Statistics}

\begin{table}[h]
\centering
\caption{Reformulation Slope Statistics}
\label{tab:slope_stats}
\begin{tabular}{lll}
\toprule
Metric & Value & Interpretation \\
\midrule
Mean slope & $-0.0214$ & Average decrease in reformulation rate per turn \\
Median slope & $0.0000$ & Half of conversations show no slope \\
Std deviation & $0.0870$ & High variance in slopes across conversations \\
$t$-statistic & $-2.310$ & Test statistic for one-sample $t$-test \\
$p$-value & $0.012$ & Statistically significant at $\alpha = 0.05$ \\
Cohen's $d$ & $-0.246$ & Small effect size \\
Sample size & 88 conversations & Conversations with $\geq 5$ turns and detected reformulations \\
\bottomrule
\end{tabular}
\end{table}

\textbf{Interpretation}: The negative mean slope ($-0.0214$) indicates that on average, reformulation rate decreases by approximately 2.1 percentage points per turn. For a 5-turn conversation, this corresponds to a $\sim$10 percentage point reduction in reformulation probability from turn 2 to turn 5. The effect is statistically significant ($p = 0.012 < 0.05$) despite small magnitude (Cohen's $d = -0.246$, below medium effect threshold of 0.5).

\subsubsection{Distribution Analysis}

The slope distribution reveals important heterogeneity:

\begin{itemize}
\item \textbf{Negative slopes}: 15 of 88 conversations (17\%) show negative slopes, indicating learning.
\item \textbf{Zero/flat slopes}: 73 of 88 conversations (83\%) show near-zero slopes (median = 0).
\item \textbf{Positive slopes}: None observed in this sample.
\end{itemize}

This bimodal distribution suggests learning is not universal. A subset of conversations exhibit clear learning curves (reformulation decreases), while the majority show flat patterns (reformulation rate constant or insufficient variation). The negative mean despite zero median indicates the learning subset has strong enough slopes to shift the population average.

\textbf{Competing explanations for median = 0}:

\begin{enumerate}
\item \textbf{Individual differences}: Some users learn AI capabilities (negative slopes) while others do not, reflecting variation in learning aptitude or engagement.
\item \textbf{Engagement decay}: Zero slopes may represent disengaged users who stop reformulating not because they learned but because they gave up. This would be indistinguishable from no learning in our data.
\item \textbf{Threshold effect}: The $\geq 5$ turn filter may select conversations already past the initial learning phase (turns 1--3), leaving only post-learning conversations with flat slopes.
\end{enumerate}

We cannot distinguish among these explanations with current data. Future work testing whether reformulation slope predicts task success (planned but not executed in this study) would differentiate learning (negative slope predicts success) from disengagement (negative slope does not predict or negatively predicts success).

\subsection{Hypothesis 2: AI Responsiveness via Diversity Correlation}

\textbf{Research question}: Does AI response diversity correlate with user query diversity?

\textbf{Finding}: Yes. Response diversity positively correlates with query diversity, indicating AI tracks query patterns. However, correlation falls below initially targeted threshold.

\subsubsection{Correlation Statistics}

\begin{table}[h]
\centering
\caption{Diversity Correlation Statistics}
\label{tab:corr_stats}
\begin{tabular}{lll}
\toprule
Metric & Value & Interpretation \\
\midrule
Pearson $r$ & $0.396$ & Medium-strength positive correlation \\
$p$-value & $< 0.001$ & Highly statistically significant \\
95\% CI & $[0.392, 0.400]$ & Confidence interval excludes zero \\
Target threshold & $r > 0.4$ & Not met (threshold failure) \\
Sample size & 169,352 conversations & All conversations with $\geq 5$ turns \\
\bottomrule
\end{tabular}
\end{table}

\textbf{Interpretation}: The correlation $r = 0.396$ indicates that approximately 15.7\% of variance in response diversity is explained by query diversity ($R^2 = 0.157$). This is a statistically robust finding ($p < 0.001$) across a large sample, but falls short of our preregistered threshold ($r > 0.4$, medium-large effect). The 95\% confidence interval $[0.392, 0.400]$ is narrow and excludes the target threshold, confirming the effect is real but weaker than predicted.

\textbf{Why below threshold?} We hypothesize Distinct-1 (lexical diversity) underestimates true responsiveness because it only captures vocabulary variation, not semantic variation. Two responses with different words but similar meanings yield high Distinct-1 but low semantic diversity. Recomputing with semantic diversity metrics (SBERT embedding variance) may yield stronger correlations, but this remains untested.

\subsubsection{Temporal Dynamics: Responsiveness Strengthens with Conversation Length}

Stratifying by conversation length reveals a striking pattern:

\begin{table}[h]
\centering
\caption{Diversity Correlation by Conversation Length}
\label{tab:temporal_corr}
\begin{tabular}{llll}
\toprule
Turn count & Pearson $r$ & Sample size & Interpretation \\
\midrule
2--3 turns & $0.04$ & 58,321 & Weak correlation in short conversations \\
4--5 turns & $0.12$ & 67,842 & Moderate correlation emerges \\
6+ turns & $0.19$ & 43,189 & Stronger correlation in longer conversations \\
\bottomrule
\end{tabular}
\end{table}

\textbf{Interpretation}: Responsiveness is not static---it strengthens as conversations extend. Short conversations (2--3 turns) show minimal correlation ($r = 0.04$), while longer conversations (6+ turns) show nearly $5\times$ stronger correlation ($r = 0.19$). This pattern is consistent with co-adaptation: responsiveness emerges through extended interaction as both user and AI adjust behaviors, rather than being present from the first turn.

\textbf{Alternative explanation}: This could reflect a confound rather than true temporal dynamics. Longer conversations accumulate more vocabulary variation (both queries and responses use more unique words over turns), artificially inflating Distinct-1 and thus correlation. Partial correlation controlling for conversation length would isolate responsiveness from length effects, but this analysis is not yet performed.

\subsection{Unexpected Finding: Heterogeneity in Learning Patterns}

The median slope = 0 finding was unexpected. Our hypothesis predicted universal learning (all users reformulate less over turns), but results show learning is conditional rather than universal. This has important implications for the bidirectional alignment framework:

\textbf{What we expected}: All conversations show negative slopes (learning), with variation only in magnitude (some users learn faster than others).

\textbf{What we observed}: Bimodal distribution where most conversations show flat slopes (no detectable learning) while a minority show strong negative slopes (clear learning).

\textbf{Theoretical revision}: Learning may be conditional on task success or user engagement. Users who successfully accomplish their goals may exhibit learning curves, while users who struggle or disengage may show flat slopes. This suggests reformulation slope alone is insufficient to measure user learning---we must account for conversation outcome (success vs abandonment) to distinguish learning from disengagement.

Future work will stratify analyses by task success and test whether negative slopes predict success (learning interpretation) or failure (disengagement interpretation). If learning interpretation holds, we expect successful conversations to have steeper negative slopes than unsuccessful conversations.

\subsection{Summary of Evidence}

Our results provide partial validation of the bidirectional alignment framework:

\textbf{Supported claims}:
\begin{itemize}
\item[$\checkmark$] User learning exists: Reformulation slope is significantly negative ($p=0.012$), validating learning curve hypothesis.
\item[$\checkmark$] AI responsiveness exists: Diversity correlation is significantly positive ($r=0.396$, $p<0.001$), validating responsiveness hypothesis.
\item[$\checkmark$] Temporal co-adaptation: Correlation strengthens with conversation length ($r=0.04 \rightarrow 0.19$), consistent with alignment emerging over interaction.
\end{itemize}

\textbf{Refuted claims}:
\begin{itemize}
\item[$\times$] Strong responsiveness threshold: Correlation $r=0.396$ falls below $r>0.4$ target, indicating weaker effect than predicted.
\end{itemize}

\textbf{Uncertain claims}:
\begin{itemize}
\item[?] Universal learning: Median slope = 0 challenges universal learning assumption. Learning may be conditional on engagement or success.
\item[?] Coupling hypothesis: Whether correlation between reformulation slope and diversity correlation predicts helpfulness beyond individual metrics remains untested (planned for future work).
\end{itemize}

Effect sizes are smaller than targeted ($d=-0.246$, $r=0.396$), but statistical significance is robust across large samples. The next step is testing whether these signals combine (coupling) to predict conversation quality, validating the full bidirectional alignment framework.
